\documentclass[11pt]{article}
\usepackage[margin=1in]{geometry}
\usepackage{booktabs}
\usepackage{hyperref}
\usepackage{amsmath}
\usepackage{url}

\title{Memo 03: MP1 Mass-Conservation Check and Updated Diagnosis}
\author{FVCOM-MP Delaware Test Cases}
\date{April 30, 2026}

\begin{document}
\sloppy
\maketitle

\section*{Purpose}
This memo updates the MP1 mass-conservation diagnosis after rerunning case
\texttt{a} with \texttt{SEDIMENT\_MODEL=T}.  The central comparison is now:
case \texttt{a} as no-floc plastic with sediment runtime on, versus case
\texttt{b} as floc-enabled but effectively null-kinetic/split-disaggregated
plastic.

\section*{Archived Analysis Outputs}
Three mass-conservation snapshots are now preserved under
\texttt{FVCOM\_MP\_test\_run/PYTHON/output/mass\_conservation/archive/}.

\begin{table}[h]
\centering
\small
\begin{tabular}{lp{0.34\linewidth}p{0.43\linewidth}}
\toprule
Version & Archive label & Meaning \\
\midrule
\texttt{v001} & \texttt{v001\_before\_case\_a\_sediment\_on} &
old comparison with case \texttt{a} using \texttt{SEDIMENT\_MODEL=F} \\
\texttt{v002} & \texttt{v002\_case\_a\_sediment\_on\_full10d} &
comparison with case \texttt{a} using \texttt{SEDIMENT\_MODEL=T} and the
original sediment BBL settings \\
\texttt{v003} & \texttt{v003\_case\_a\_sediment\_on\_no\_bbl} &
case-\texttt{a} diagnostic with \texttt{SEDIMENT\_MODEL=T} but sediment BBL
roughness switches disabled \\
\bottomrule
\end{tabular}
\end{table}

The live analysis folder has also been regenerated from the \texttt{v003}
state.  The archived files include \texttt{mass\_summary.csv},
\texttt{mass\_timeseries\_long.csv}, and all current comparison plots.

\section*{Current Evidence}
All three current MAT files cover the same FVCOM model-day window,
\(58485.0\)--\(58494.957031\), or comparison day \(0.00\)--\(9.96\).

\begin{table}[h]
\centering
\begin{tabular}{lrrrrr}
\toprule
Case & Sediment runtime & Final water & Final bed & Final total & Retained vs. \texttt{a} \\
 & & (kg) & (kg) & (kg) & (\%) \\
\midrule
\texttt{a} & on, BBL off & 573.28 & 0.00 & 573.28 & 100.0 \\
\texttt{b} & on & 371.85 & 0.00 & 371.85 & 64.9 \\
\texttt{c} & on & 504.26 & 24.29 & 528.55 & 92.2 \\
\bottomrule
\end{tabular}
\end{table}

The major finding from \texttt{v002} remains that enabling sediment runtime in
case \texttt{a} reduced the final domain MP1 mass from the old archived value
\(1001.25~\mathrm{kg}\) to \(582.41~\mathrm{kg}\).  The new \texttt{v003}
BBL-off diagnostic ends at \(573.28~\mathrm{kg}\), only
\(9.13~\mathrm{kg}\) below the \texttt{v002} BBL-on value.  Therefore a large
part of the old case-\texttt{a}/case-\texttt{b} gap was not caused by the split
floc path alone, and the BBL-off test suggests the sediment-runtime effect is
not primarily controlled by sediment BBL roughness.

The remaining null-floc mismatch is still meaningful.  With sediment runtime
matched in the latest \texttt{v003} comparison, case \texttt{b} ends
\(201.43~\mathrm{kg}\) below case \texttt{a}, retaining \(64.9\%\) of the
sediment-on no-floc baseline.  Case \texttt{c} ends
\(44.73~\mathrm{kg}\) below case \texttt{a}, retaining \(92.2\%\), with
\(24.29~\mathrm{kg}\) stored in the bed.

\section*{New Negative Test: BBL-Off Case \texttt{a}}
To test whether sediment affected plastic only through sediment-controlled
bottom roughness, a case-\texttt{a} diagnostic sediment input was prepared:
\texttt{generic\_sediment\_a\_no\_bbl.inp}.  This keeps
\texttt{SEDIMENT\_MODEL=T} and suspended sediment transport enabled, but sets
\texttt{MB\_BBL\_USE=F}, \texttt{SG\_BBL\_USE=F},
\texttt{SSW\_BBL\_USE=F}, and \texttt{SSW\_CALC\_ZNOT=F}.

The resulting analysis is archived as
\texttt{v003\_case\_a\_sediment\_on\_no\_bbl}.  It showed almost no change
relative to the sediment-on case-\texttt{a} run with the original BBL settings:
case \texttt{a} ended at \(573.28~\mathrm{kg}\) in \texttt{v003}, compared
with \(582.41~\mathrm{kg}\) in \texttt{v002}.  This is important negative
evidence: the sediment/plastic interference is unlikely to be explained only by
the Sherwood/Signell/Warner BBL roughness feedback.  It also means that either
the HPC build did not actively use that BBL path, or another sediment-enabled
code path dominates the plastic difference.

\section*{New Isolation Test: Delayed \texttt{Advance\_Sed}}
A sharper test was then prepared using
\texttt{generic\_sediment\_a\_inactive.inp}.  This keeps
\texttt{SEDIMENT\_MODEL=T} and the sediment setup path active, but sets
\texttt{SED\_START = 999999999}.  In \texttt{mod\_sed.F:883},
\texttt{Advance\_Sed} returns immediately when
\texttt{iint < sed\_start}.  Therefore this test keeps sediment initialization,
the sediment/plastic compile path, and the river-forcing optional-argument path,
but removes active runtime sediment transport.

The recent partial run is already diagnostically useful: the
\texttt{total\_mass\_by\_case.png} curve is reported to overlay the original
\texttt{v001} behavior.  Assuming the completed run preserves this result, the
main interference is not from sediment setup, sediment namelist presence, or the
compile-time river-forcing path.  It is instead localized to code executed after
the early return in \texttt{Advance\_Sed}, i.e. the active sediment operator.

\section*{Cap, Bed, and Split Checks}
The hard concentration cap remains inactive in the present outputs.  The
maximum extracted concentrations are far below the code cap of
\(100~\mathrm{kg\,m^{-3}}\):
\[
\max C_a = 2.41\times 10^{-4}, \qquad
\max C_b = 6.98\times 10^{-3}, \qquad
\max C_c = 1.45\times 10^{-2}~\mathrm{kg\,m^{-3}}.
\]
The cap-threshold counters are zero for all cases with available counters.
Thus \texttt{mod\_plast.F:1242--1243} is still a non-conservative code site,
but it is not active in this 10-day test.

For the null-floc comparison, bed exchange is also not the leading cause.
Cases \texttt{a} and \texttt{b} both end with zero sampled bed MP1 mass and
zero sampled deposition/erosion exchange.  Case \texttt{c} does exchange with
the bed, so it should be interpreted separately.

The saved floc split fields are internally consistent:
\[
\texttt{mp1\_agg}+\texttt{mp1\_dis}\approx\texttt{mp1}.
\]
The maximum total split-minus-combined mismatch is zero for case \texttt{b}
and \(1.83\times10^{-7}~\mathrm{kg}\) for case \texttt{c}.  Saved-output
recombination is therefore not the present sink.

\section*{Sediment/Plastic Call Trace}
The setup and timestep order make the sediment-runtime finding plausible.

\begin{enumerate}
\item \textbf{Setup order.} In \texttt{fvcom.F:643--651},
\texttt{SETUP\_SED} is called before \texttt{SETUP\_PLAST}.  Thus the
sediment module is initialized first when \texttt{SEDIMENT\_MODEL=T}.

\item \textbf{Timestep order.} In \texttt{internal\_step.F:1092--1137},
\texttt{ADVANCE\_SED} is called before \texttt{ADVANCE\_PLAST}.  The earlier
offline-style block in \texttt{internal\_step.F:213--235} has the same
ordering.  Plastic therefore sees the post-sediment state each internal step.

\item \textbf{Sediment dynamics.} \texttt{mod\_sed.F:859--1156} advances
active-layer thickness, settling/deposition, erosion, bed stratigraphy/bottom
properties, horizontal advection, FCT limiting, vertical diffusion, OBC/source
boundary conditions, nonnegative/upper clipping, and sediment exchange.

\item \textbf{Plastic dynamics.} \texttt{mod\_plast.F:914--1851} computes
ambient density and floc rate fields, updates plastic active-layer quantities,
deposition/erosion/update-bottom, horizontal advection, vertical diffusion,
kinetics, OBC/source boundary conditions, clipping, MPI exchange, and final
recombination.

\item \textbf{Floc coupling to sediment.} In
\texttt{mod\_plast.F:4619--4815}, the aggregation rate calculation reads
\texttt{sed(ised)\%conc}, \texttt{sed(ised)\%Sd50}, and
\texttt{sed(ised)\%Wset}.  Thus active floc interaction is explicitly tied to
sediment size and suspended sediment concentration.  In case \texttt{b}, the
saved output still has \texttt{mp1\_agg}=0, so the current \texttt{b} run is
effectively a split-disaggregated/null-floc test despite executing the floc
branch.
\end{enumerate}

\section*{Diagnosis}
\begin{enumerate}
\item \textbf{Sediment runtime changed domain retention.} The drop from
\(1001.25~\mathrm{kg}\) in \texttt{v001} to \(582.41~\mathrm{kg}\) in
\texttt{v002} shows that turning on the sediment model changes the transport
environment and/or open-boundary export enough to strongly affect domain MP1
mass.  This is not necessarily a plastic mass sink by itself; the plotted
domain mass does not include cumulative open-boundary export.

\item \textbf{BBL roughness is now demoted.} The archived \texttt{v003}
\texttt{generic\_sediment\_a\_no\_bbl.inp} test changed case-\texttt{a} final
mass by only \(9.13~\mathrm{kg}\) relative to the \texttt{v002} BBL-on
sediment case.  Bottom-stress feedback should remain documented, but it is no
longer the leading explanation for the present discrepancy.

\item \textbf{The no-floc plastic path itself changes when sediment is
compiled/enabled.} In \texttt{mod\_plast.F}, the update enters a
\texttt{\#if defined(SEDIMENT)} block before checking
\texttt{floc\_intract}.  Even with \texttt{PLAST\_FLOC=F}, case \texttt{a}
uses the sediment-build no-floc branch around
\texttt{mod\_plast.F:1248--1278}.  Therefore an old case-\texttt{a} run built
without \texttt{-DSEDIMENT} is not a strict code-path match to a sediment-build
case-\texttt{a} run.

\item \textbf{River forcing/source updating is now demoted.} The river update
call still changes at compile time: with plastic only,
\texttt{bcond\_gcn.F:645} calls \texttt{UPDATE\_RIVERS} with
\texttt{PLAST=PLASTDIS}; with sediment and plastic,
\texttt{bcond\_gcn.F:648} calls it with both \texttt{SED=SEDDIS} and
\texttt{PLAST=PLASTDIS}.  However, the delayed-\texttt{Advance\_Sed} test keeps
this compile/setup path while recovering the \texttt{v001}-like mass curve.
Thus river optional-argument handling is unlikely to be the primary cause of
the present case-\texttt{a} discrepancy, although cumulative river MP input
should still be logged for closure.

\item \textbf{The leading suspect is now active \texttt{Advance\_Sed}
side effects.} \texttt{Advance\_Sed} is called before
\texttt{Advance\_Plast} each internal step.  With
\texttt{MORPHO\_MODEL=F}, bathymetry should not be updated, and with BBL
disabled the roughness route is weakened.  The delayed-\texttt{Advance\_Sed}
test localizes the main discrepancy to the active sediment operator, especially
the shared scalar-transport calls and any shared model state they leave behind
before plastic advances.

\item \textbf{Shared scalar-transport state is a concrete mechanism to test.}
Inside \texttt{mod\_sed.F:1016--1028}, sediment calls
\texttt{adv\_scal\_backward} and \texttt{fct\_sed}.  The shared scalar routine
writes global bottom-gradient arrays \texttt{PFPXB}/\texttt{PFPYB} in
\texttt{mod\_scal.F:911--912}.  No-floc plastic should not directly read
sediment concentration, so a case-\texttt{a} response to active sediment
transport points toward shared-state contamination or another common scalar
side effect rather than physical floc aggregation.

\item \textbf{A sediment exchange mismatch can affect floc cases.} After
sediment clipping, \texttt{mod\_sed.F:1099--1104} updates
\texttt{sed(i)\%conc}; the following MPI exchange at
\texttt{mod\_sed.F:1117--1122} exchanges \texttt{sed(i)\%cnew}, not the clipped
\texttt{sed(i)\%conc}.  Floc-rate calculations later read
\texttt{sed(ised)\%conc}.  This is less likely to explain the no-floc
case-\texttt{a} effect by itself, but it is a strong candidate for b/c
rank-boundary oscillations and should be corrected or instrumented.

\item \textbf{The remaining case-\texttt{a}/case-\texttt{b} gap is still a
split-path problem after sediment state is matched.} With
\texttt{SEDIMENT\_MODEL=T} in both cases, the remaining
\(210.56~\mathrm{kg}\) deficit in case \texttt{b} points back to differences
between combined \texttt{conc} and split \texttt{conc\_d} operators.

\item \textbf{The leading code-level suspects are water-column transport and
boundary handling.} With cap counts zero and no sampled bed storage in
case \texttt{b}, the likely divergence occurs during split
\texttt{adv\_scal}, split \texttt{vdif\_scal}, split
\texttt{bcond\_scal\_OBC}, split \texttt{bcond\_scal\_PTsource}, or
nonnegative clipping before output.

\item \textbf{The split branch still has one awkward exchange point.} Before
OBC handling, the floc branch exchanges \texttt{cnew\_a},
\texttt{cnew\_d}, and combined \texttt{cnew}
(\texttt{mod\_plast.F:1513--1516}).  At that point the meaningful prognostic
state is split; combined \texttt{cnew} is not rebuilt until
\texttt{mod\_plast.F:1715--1716}.  This may be harmless, but it is not
strictly symmetric with the no-floc path.

\item \textbf{Final recombination improved the budget but did not solve it.}
The patched final order now clips and exchanges \texttt{cnew\_a}/
\texttt{cnew\_d}, then rebuilds
\(\texttt{conc}=\texttt{cnew}=\texttt{conc\_a}+\texttt{conc\_d}\).  The mass
recovery after that change confirms it addressed a real asymmetry, but the
remaining deficit shows another divergence exists upstream.

\item \textbf{Direct floc-rate reads from sediment should be bypassed in a
true null-kinetics test.} The aggregation/disaggregation routines read
\texttt{sed(ised)\%conc}, \texttt{Sd50}, and \texttt{Wset}.  For a strict
``kinetics off'' symmetry test, \(\lambda_a=\lambda_d=0\) should be imposed
without calling sediment-dependent rate calculations at all.  That removes any
possible sediment-size or sediment-concentration dependence from case
\texttt{b}.

\item \textbf{Case \texttt{c} is not a pure conservation comparator.} Case
\texttt{c} retains much more total mass than case \texttt{b}, but it has active
split components and bed storage.  It is useful for full floc behavior, while
case \texttt{b} remains the sharper test for whether split \texttt{conc\_d}
can reproduce combined \texttt{conc}.
\end{enumerate}

\section*{Focused Next Checks}
\begin{enumerate}
\item Use \texttt{v003} as the latest pre-inactive-\texttt{Advance\_Sed}
archive for future plots and memo updates.  If the delayed-\texttt{Advance\_Sed}
run finishes with the same \texttt{v001}-like behavior, archive that analysis as
\texttt{v004} and record it as the decisive ``sediment initialized but runtime
operator inactive'' comparison.

\item Bisect \texttt{Advance\_Sed} using temporary early returns.  First return
after \texttt{calc\_active\_layer}; then after
\texttt{calc\_deposition}/\texttt{calc\_erosion}/bed updates but before
\texttt{adv\_scal\_backward}; then after
\texttt{adv\_scal\_backward}/\texttt{fct\_sed}; then after
\texttt{vdif\_scal}/MPI exchange but before OBC/source/clipping.  This should
pinpoint which block first changes plastic-domain mass.

\item Add a targeted shared-state checksum immediately before and after
\texttt{Advance\_Sed}, before \texttt{Advance\_Plast}.  Track small summaries
for \texttt{PFPXB}, \texttt{PFPYB}, \texttt{D}, \texttt{DT}, \texttt{DZ},
\texttt{KH}, \texttt{U}, \texttt{V}, \texttt{UARD\_OBCN}, and the shear-stress
array passed to plastic.  Any non-sediment state changed by
\texttt{Advance\_Sed} becomes an immediate smoking gun.

\item Run a narrow save/restore test around sediment \texttt{adv\_scal\_backward}
for \texttt{PFPXB}/\texttt{PFPYB}.  If restoring these arrays removes the
case-\texttt{a} sediment-runtime effect, the coupling mechanism is shared scalar
gradient state.

\item Test \texttt{SED\_ONED=T} as an intermediate isolation.  This skips the
horizontal sediment advection/OBC/source path in \texttt{Advance\_Sed} while
leaving more of the vertical/bed sediment machinery active.  A return toward
\texttt{v001} would implicate the horizontal scalar transport block.

\item Add a river-source budget log for MP1.  At each report interval, record
instantaneous \(\sum_r Q_r C_{\mathrm{plast},r}\) and cumulative
\(\sum_t\sum_r Q_r C_{\mathrm{plast},r}\Delta t\).  This will determine
whether the apparent domain-mass loss is actually lower source input or
stronger export.

\item Add a short internal checkpoint diagnostic for one small matched run.
Record domain-integrated mass for \texttt{conc}, \texttt{conc\_d}, and
\(\texttt{conc\_a}+\texttt{conc\_d}\) after: start, deposition, erosion/update
bottom, \texttt{adv\_scal}, \texttt{vdif\_scal}, OBC, point-source BC,
nonnegative clamp, MPI exchange, and recombination.

\item For the null-floc case, add a temporary switch or guard so
\(\lambda_a=\lambda_d=0\) is assigned directly and the sediment-dependent floc
rate routines are not called.

\item For b/c, reconcile sediment post-clipping exchange: exchange
\texttt{sed(i)\%conc} after clipping, or set \texttt{sed(i)\%cnew =
sed(i)\%conc} before exchanging.  This should be treated as a diagnostic fix
first, because it affects the sediment state read by the floc-rate routines.

\item Explicitly accumulate mass removed by nonnegative clipping.  The hard cap
should remain zero; any nonzero clipping loss would be immediately
informative.

\item If serial and parallel runs differ, focus on the two exchange locations:
\texttt{mod\_plast.F:1513--1516} before OBC and
\texttt{mod\_plast.F:1677--1716} at final recombination.
\end{enumerate}

\end{document}
